% ECONOMICS_PROBLEM: {"id": "I32-241", "title": "Identifying genuine poverty traps", "jel_code": "I32", "source_id": "OP-0370"}
% !TeX program = xelatex
\documentclass[11pt,a4paper]{article}
\usepackage[margin=23mm]{geometry}
\usepackage{fontspec}
\setmainfont{Latin Modern Roman}
\usepackage{amsmath,amssymb,mathtools}
\usepackage{xcolor,enumitem,booktabs,longtable,array,xurl,hyperref}
\definecolor{muted}{HTML}{53636C}
\hypersetup{colorlinks=true,urlcolor=blue,linkcolor=blue}
\setlength{\parindent}{0pt}
\setlength{\parskip}{5pt}
\newcommand{\R}{\mathbb R}
\newcommand{\E}{\mathbb E}
\newcommand{\N}{\mathbb N}
\newcommand{\ind}{\mathbf 1}
\newcommand{\Var}{\operatorname{Var}}
\newcommand{\Cov}{\operatorname{Cov}}
\newcommand{\supp}{\operatorname{supp}}
\DeclareMathOperator*{\argmin}{arg\,min}
\DeclareMathOperator*{\argmax}{arg\,max}
\newcommand{\entrymeta}[3]{{\sffamily\small\color{muted}\textbf{\texttt{#1}}\quad|\quad #2\par #3\par}}
\newcommand{\Problemref}[1]{\texttt{#1}}
\newcommand{\JELref}[2]{\texttt{#2} (JEL \texttt{#1})}
\begin{document}
\section*{Identifying genuine poverty traps}
\textbf{Problem ID:} \texttt{I32-241}\quad \textbf{JEL:} \texttt{I32}\par
% BEGIN ECONOMICS STATEMENT
\label{problem:OP-0370}
{\sffamily\small\color{muted}Original subject group: Development and poverty\par}
\label{definitions:problem:30007652}
\textbf{Audit classification:} Background; proposed extension. Authors and NBER identifier checked. The critique challenges broad threshold identification rather than certifying a universal open conjecture; deterministic/stochastic model discrimination is editorial.
\subsection*{Definitions and precise research target}
\textbf{Complete editorial problem.} Fix a population of households, an institutional environment, and a horizon long enough to distinguish recovery from transitory dynamics. Household \(i\)'s productive assets follow \(a_{i,t+1}=f(a_{it},\theta_i,\eta_{it})\), where \(\theta_i\) is persistent heterogeneity and \(\eta_{it}\) a shock with specified distribution. For a type \(\theta\), define expected drift \(d(a,\theta)=E[f(a,\theta,\eta)-a\mid\theta]\). An asset-threshold trap in this particular model requires at least two attracting long-run regimes and an intervening unstable threshold, rather than simply low assets or slow growth. For the deterministic benchmark, specify attracting fixed points and their basins. When shocks permit crossings, replace absorbing-regime language by a declared Markov-kernel hitting-time or metastability criterion; expected drift alone does not establish multiple attracting regimes. The identification task is to distinguish such type-specific nonconvex dynamics from a mixture of households with different smooth dynamics, measurement error, and temporary shocks. Use randomized asset increments over adequate support, repeated measurements, and exogenous shocks to identify or bound \(f\) and long-run responses. Determine whether any common population threshold is empirically defensible. Do not impose it because a fitted aggregate curve appears nonlinear. The cited methodological critique questions whether broad threshold evidence is attainable under heterogeneity, so the task is scoped identification and model discrimination rather than an asserted universal poverty-trap conjecture.

\textbf{Definitions and admissible scope.} A precise deterministic benchmark is \(a_{t+1}=f_\theta(a_t)\) on a compact nonnegative asset interval, with a continuous differentiable map and three fixed points \(a_L<a_U<a_H\). Require \(|f'_\theta(a_L)|<1\), \(|f'_\theta(a_H)|<1\), and \(|f'_\theta(a_U)|>1\), plus stated basins of attraction. This is a model definition of a threshold trap, not an empirical conclusion. A noisy transition is instead a Markov kernel \(K_\theta(a,da')\). If shocks connect all assets, two absorbing long-run classes may not exist; define metastability through expected hitting times or finite-horizon transition probabilities. An aggregate conditional drift curve from a mixture of types need not equal any type's map. The identification problem therefore concerns kernels, heterogeneity, support, and measurement rather than declaring multiple stable regimes from one nonlinear fitted curve. Define the stochastic hitting time \(\tau_H=\inf\{t\ge0:a_t\in H\}\), with \(\inf\varnothing=+\infty\), for a declared measurable asset set \(H\). A finite-horizon recovery target is \(h_T(a,\theta)=\Pr_\theta(\tau_H\le T\mid a_0=a)\); a long-run escape target can be \(E_\theta\tau_H\), permitting infinity. The model class specifies the Markov transition, persistent-type distribution, shock serial dependence or Markov augmentation, and measurement-error law. Matching finite observations constrains these objects but need not identify infinite-horizon behavior. State separately whether the requested output is a deterministic threshold, finite-horizon kernel bounds, or a stochastic escape-time comparison.
\subsection*{Variants and assumption changes}
\begin{enumerate}
\item Randomized increments (editorial): observe several transfer sizes covering both sides of a hypothesized threshold and repeated assets. Identify conditional transition kernels or bounds by baseline type under an explicit measurement and uptake model.
\item Stochastic persistence (editorial): fix asset sets \(L,H\), horizon \(T\), and the probability \(P_\theta(\tau_H\le T\mid a_0=a)\). Compare smooth single-regime kernels with kernels having long expected escape times; a finite experiment does not by itself establish infinite-horizon absorbing traps.
\end{enumerate}
\subsection*{References and source audit}
\textbf{Support and access.} Authors and NBER identifier checked. The critique challenges broad threshold identification rather than certifying a universal open conjecture; deterministic/stochastic model discrimination is editorial. \textbf{Locator:} IPR WP-26-24 abstract, final paragraph.
\begin{enumerate}\raggedright
\item\hypertarget{definitions:source:30007652:1}{} Dean Karlan; Amol Singh Raswan; Christopher Udry. 2026. \emph{The Sisyphean Pursuit of Evidence for Poverty Traps}. Abstract, final paragraph.
\url{https://www.ipr.northwestern.edu/our-work/working-papers/2026/wp-26-24.html}.
DOI: \url{https://doi.org/10.3386/w35253}.
\end{enumerate}

\subsection*{Common conventions: Source mathematical conventions}

These conventions apply to each entry unless explicitly replaced.

\textbf{Models and identification.} A primitive model \(m\) specifies all probability laws, feasible choices, information, equilibrium and measurement rules used by its target. The admissible class \(\mathcal M\) is fixed independently of outcomes. If \(P_O(m)\) is its observable probability law and \(\tau(m)\) the target, its identified set at observable law \(P\) is
\[
 \mathcal I_\tau(P)=\{\tau(m):m\in\mathcal M,\ P_O(m)=P\}.
\]
Point identification means that this set is a singleton. An empty set signals incompatibility of data and assumptions. Coordinate bounds need not describe the joint set. Bounds are sharp only if every retained target is attainable or arbitrarily approachable by compatible models. Finite bounds require explicit restrictions on unobserved quantities; unrestricted confounding, hidden wealth, extrapolation or arbitrary equilibrium selection can yield uninformative sets.

\textbf{Causal interventions.} A potential outcome \(Y_i(p)\) is the outcome under a complete policy regime \(p\), including timing, financing, information and market spillovers. The target population and units are fixed before treatment. With interference, use the complete assignment vector or a justified exposure mapping; individual no-interference notation alone cannot identify an economy-wide intervention. A difference of mean potential outcomes does not require the cross-world joint distribution. Conditioning on post-treatment migration, survival, enrollment or employment changes the estimand and needs separate justification.

\textbf{Statistical statements.} An estimator, confidence set, forecast or test specifies a sampling law, model class and loss/coverage criterion. Uniform \((1-\alpha)\)-coverage means \(\inf_{m\in\mathcal M}\Pr_m\{\tau(m)\in C_n\}\geq1-\alpha\), or an explicitly asymptotic version. The whole parameter space trivially has coverage, so informativeness requires a stated width, risk or power objective. A forecasting improvement is relative to a named benchmark, information set and evaluation population.

\textbf{Equilibrium and optimization.} An equilibrium comprises feasible plans, optimality, beliefs where needed, budget constraints and all market/resource clearing. The primitives or a stated benchmark define each correspondence \(\mathcal E(p;m)\). Nonemptiness is assumed or proved, never inferred just from compact policy sets. A selected-equilibrium target uses a declared selection rule; a robust target quantifies over all permitted equilibria. Use suprema or infima unless attainment is established. An \(\varepsilon\)-optimal policy is feasible and within \(\varepsilon\) of the finite optimum value. Report an empty feasible set separately.

\textbf{Welfare and accounting.} Normative utility, interpersonal normalization, social weights, numeraire, discounting and the use of public receipts are primitives. Behavioral utility need not coincide with the welfare criterion. Household welfare includes ownership income and rents once; adding profits again to compensating variation generally double-counts them. A monetary equivalent is defined by an explicit transfer and price convention, with monotonicity and range conditions. Average effects, mechanism contributions, transfers and welfare losses are different targets. Infinite sums require convergence and dynamic feasible plans require terminal or no-Ponzi conditions.

\textbf{Source versions.} A working-paper date, revision date and journal publication year are distinguished. A DOI for a journal version is not automatically the DOI of an earlier manuscript. Locators refer to the stated version and printed pages where specified. A metadata-only check does not verify a claim about a full-text conclusion. Every entry's source subsection narrows the provenance claim accordingly.
% END ECONOMICS STATEMENT
\end{document}
